\section{Introduction}
\label{sec:intro}

\subsection{Grid access as the scarce resource}

Electricity demand from data centres is growing faster than at any time in the last two decades, driven by the training and serving of artificial-intelligence models \cite{iea2025,masanet2020}. In most mature markets the limiting factor is no longer the price or availability of energy but the right to draw power from the grid at a given node. Spain is a clear example. By 1 October 2026 the transmission system operator (TSO), Red El\'ectrica, had granted 18.5\,GW of demand access on the transmission grid, of which 6.3\,GW corresponded to data centres, none of them yet in service \cite{ree2026mapa}; only about one in four transmission nodes can still accept new demand \cite{ree2026nudos}. The Government has reacted with capacity tenders at saturated nodes and an annual reservation fee \cite{rdl72026}, a new category of flexible access permits \cite{cnmc2026flex}, and a draft royal decree that conditions the access permit of large data centres on efficiency and hourly renewable-matching requirements \cite{prd2026,ashurst2026}.

At the same time, part of the access capacity already granted at transmission level is structurally underused. Railway infrastructure holds 3.55\,GW of granted demand access on the Spanish transmission grid \cite{ree2026mapa}. High-speed traction substations are connected at 220 or 400\,kV and are sized for the worst minutes of the timetable under the loss of an adjacent substation, yet a high-speed line is, electrically, a sparse sequence of short power pulses: the mean load of a substation is a small fraction of its peak. Because the methodology used to compute node margins subtracts the full capacity of every granted permit \cite{cnmc2025capacidad}, this headroom is invisible to other consumers. It is \emph{capacity}, not energy, that is idle: regenerative braking returned to the transmission grid is a minor flow in 25\,kV AC high-speed systems, and this paper does not rely on it.

\subsection{Related work}

Three bodies of literature meet in this problem.

\emph{Data centres as flexible loads.} The potential of data centres for demand response has been studied for more than a decade, from early surveys and pricing mechanisms \cite{wierman2014,liu2014} to spatial load shifting between sites \cite{zheng2020,zhang2020}. Recent work changes the scale of the argument: Norris et al.\ estimate that the United States grid could integrate on the order of 100\,GW of new large loads if these accept curtailment during a fraction of a percent of the year \cite{norris2025}, and a field demonstration on a production GPU cluster sustained a 25\,\% power reduction for three hours without violating quality of service \cite{colangelo2025}. Industry programmes now test this at several sites \cite{epri2024}. Two caveats are equally recent. Training workloads can produce fast, large power swings \cite{li2024}, and voltage-sensitive data centres can disconnect simultaneously: on 10 July 2024 about 1,500\,MW of data-centre load in Virginia transferred to on-site supply after a transmission fault \cite{nerc2025}. Flexibility also interacts with clean-energy procurement: hourly matching is far more demanding than annual matching \cite{xu2024,riepin2024}.

\emph{Electrical behaviour of AC railways.} The traction load is single-phase, pulsating and partly regenerative. Its power-quality impact on the supplying grid, mainly voltage unbalance, is well characterised \cite{gazafrudi2015}, and advanced supply schemes with static converters have been proposed to remove neutral sections and balance the load \cite{krastev2016,pilo2015}. Energy storage and regenerative-braking recovery have been studied extensively, mostly for DC urban systems \cite{gonzalezgil2013,khodaparastan2019,ratniyomchai2014}, together with wider energy-saving measures \cite{douglas2015}. Stationary storage and photovoltaic generation have been co-optimised with the loads of a railway station \cite{sengor2018}.

\emph{Sharing railway assets with the power system.} Rail corridors have been proposed as carriers of mobile storage \cite{moraski2023}. Closest to this work, a recent study estimates the non-railway load that the DC traction substations of the Rotterdam metro could host, between 2 and 5\,MW, and finds that the binding limit is sometimes the contractual grid-connection capacity rather than the equipment \cite{dwarkasing2026}. In Spain, the competition authority has recommended that the infrastructure manager publish the location and consumption of its high-speed substations and offer adjacent land for renewable generation \cite{cnmc2023adif}.

\subsection{Gap and contributions}

To the best of our knowledge, no study has quantified the data-centre load that high-speed (25\,kV AC, transmission-connected) traction substations could host by sharing their grid-access capacity, nor specified how such an installation should be built, protected and regulated. The question has become more demanding in 2026. Two serious accidents on the Madrid--Seville high-speed line, in January and October \cite{adamuz2026,urda2026}, repeated cable thefts \cite{cable2026} and a large data breach at the infrastructure manager and the incumbent operator \cite{cyber2026} have put railway safety and security at the centre of public debate and led to changes at the head of the transport ministry and of the infrastructure manager \cite{infobae2026}. None of these events involved the traction power supply, but they set the standard that any proposal for third-party use of railway electrical infrastructure must meet: it must be demonstrably unable to degrade railway operation.

The contributions of this paper are:
\begin{enumerate}
\item A reproducible, open model that derives second-by-second substation loads from public timetables and maps them onto the access capacity actually granted to each railway bay, for a high-density line (Madrid--Barcelona) and a low-density line (Madrid--Seville).
\item The hosting capacity for data centres under current rules and under a proposed \emph{shared access permit with priority}, without storage, with a curtailable share and with optimally sized batteries, including sensitivity to traffic growth, access level and the averaging window of the joint limit.
\item An implementation architecture, with its control hierarchy, cost and a staged pilot, showing that the concept can be built with standard equipment.
\item A security analysis (train priority, ride-through, systemic risk, physical and cyber security) that turns the lessons of recent railway and grid incidents into design requirements.
\item The regulatory pathway and the interaction with the Spanish data-centre decree, including the cost of hourly renewable matching with and without flexibility.
\end{enumerate}

Section~\ref{sec:methods} describes the system and the model; Section~\ref{sec:results} gives the results; Section~\ref{sec:arch} presents the architecture and its feasibility; Section~\ref{sec:security} deals with security and resilience; Section~\ref{sec:discussion} discusses regulation, limitations and future work; Section~\ref{sec:conclusions} concludes.
